% ============================================================
% Paper 2 - From Local Curvature to the Weil Functional:
%           An Explicit Construction
%
% Author:   Ulrich Tehrani
% Date:     May 2026
% Version:  10
% Zenodo:   https://doi.org/10.5281/zenodo.19106992
% License:  CC BY 4.0
%
% Changelog
% ---------
% v1--v4 (March 2026): see project history
%
% v5 (April 2026)
%   - Introduction framing softened: only the prime-side
%     construction is claimed; spectral inequality marked open.
%   - Abstract clarified: spectral inequality is RH-strength.
%   - Status remark added before the open inequality.
%
% v10 (May 2026)
%   - Reader Contract verification description corrected:
%     "Both numerical statements are verified on a finite grid
%     and are labelled [numerical] in their theorem headings"
%     applied "verified on a finite grid" to both observations,
%     but Observation 4.3 (D approx 9.471) is a finite prime-
%     truncation plus tail-bound computation, not a grid check;
%     only Observation 6.1 is a grid check. Replaced with
%     "Both numerical statements are finite computations from
%     the stated parameters and are labelled [numerical] in
%     their theorem headings". The §6 setup phrase
%     "renormalised local curvature on a finite grid in kappa"
%     is correct as section context and is left unchanged.
%   - No changes to numerical values, theorem statements, or
%     proofs. The §5 prose is unchanged.
%
% v9 (May 2026)
%   - Two consistency edits:
%     (a) Abstract "Structure" paragraph: the closing sentence
%         in v8 said "non-circular and proof-complete". Since
%         §3 contains the conditional Proposition 3.1, the
%         "proof-complete" attribution overstated the status of
%         §§2-5. Replaced by "non-circular, with proved and
%         conditional components explicitly separated". (The §8
%         Non-Circularity section was already paragraph-by-
%         paragraph specific and did not use "proof-complete"
%         phrasing; no change needed there.)
%     (b) Numerical-label discipline in Lemma 4.2: in v8 the
%         lemma "Convergence of the diagonal energy; proved"
%         contained the numerical sentence
%         "Numerically, D approx 9.471 ... tail bound below
%         0.009", mixing a proved statement with a numerical
%         observation under the same "; proved" label. Lemma
%         4.2 now contains only the convergence claim. The
%         numerical evaluation moves into a new Observation 4.3
%         "Diagonal energy value; numerical" with label
%         obs:D-value, immediately following Lemma 4.2. Reader
%         Contract updated to point at the new observation.
%         The Reader Contract no longer lists the sawtooth
%         behaviour as a numerical bullet, since the sawtooth
%         statement lives in §5 prose, not in a labelled
%         theorem environment; the §5 prose itself is unchanged.
%   - No changes to numerical values, theorem statements
%     (other than the Lemma 4.2 / new Observation 4.3 split
%     and the abstract-closing wording), or proofs.
%
% v8 (May 2026)
%   - Five edits, all in the prime-side identity area:
%     (a) Proposition 3.1 (prime-side identity): status
%         changed from 'proved' to 'conditional', with the
%         named assumption being the standard remainder
%         bookkeeping for the Weil explicit formula applied
%         to antisymmetric Schwartz test functions at fixed
%         (sigma, kappa) as epsilon -> 0. Rationale: the
%         theorem heading carried '; proved' while the proof
%         body was a 'Sketch', which is incoherent under the
%         strict series labelling discipline.
%     (b) Bookkeeping wording in the Proposition 3.1 sketch
%         corrected: the apparent self-contradiction between
%         "the archimedean contribution to W is zero" and
%         "the archimedean piece gives psi_1(sigma)" was the
%         result of conflating two distinct objects in the
%         Lagarias decomposition. The constant term
%         (proportional to hat g^*(0)^2 and vanishing by
%         antisymmetry) and the archimedean integral
%         (contributing psi_1(sigma)) are now distinguished
%         explicitly.
%     (c) Abstract reorganised: a new "What is conditional"
%         category collects the prime-side identity; "What is
%         proved" retains only the three lemmas (Admissibility,
%         f_p positivity, D convergence). Reader Contract
%         follows the same reorganisation.
%     (d) Observation 6.1 narrowed to a strictly finite-grid
%         statement: title changed from "Asymptotic stability
%         of the comparison" to "Finite-grid stability of the
%         comparison". The earlier claim that violations occur
%         "only on the small-kappa window kappa <= 25 (in the
%         prime gap between 23 and 29)" is withdrawn, since
%         the only tested point with kappa < 26 is kappa = 20,
%         which is not in [24, 28]. The observation now states
%         only that kappa = 20 is the single grid violation.
%         Surrounding text in §6 and Open Problem 7.1 updated
%         accordingly.
%     (e) Mertens-constant formula in the §5 remark corrected.
%         The v7 expression
%           gamma_M = gamma + sum_p [log(p/(p-1)) - 1/p]
%         has the wrong sign for the quoted value 0.2615.
%         Replaced by the standard Meissel-Mertens formula
%           gamma_M = gamma + sum_p [log(1 - 1/p) + 1/p] ~ 0.2615.
%   - No changes to numerical values, theorem statements
%     (other than the Proposition 3.1 status label and the
%     Observation 6.1 title and statement narrowing), or
%     proofs.
%
% v7 (April 2026)
%   - Bibliographic and positioning revision pass:
%     (a) Edwards bibitem: publisher corrected from Dover to
%         Academic Press (New York, 1974).
%     (b) Connes reference replaced: Connes-Consani 2021
%         "Weil positivity and trace formula" -> Connes 1999
%         "Trace Formula in Noncommutative Geometry and the
%         Zeros of the Riemann Zeta Function". The 1999 paper
%         is the canonical operator-theoretic reference for
%         the Hilbert-Polya/noncommutative-spectral-
%         interpretation programme that the present
%         construction must be distinguished from; Paper 6 of
%         the series cites the same source for the same
%         purpose.
%     (c) Three positioning bibitems added (Berry-Keating,
%         de Branges, Burnol), matching Paper 6's positioning
%         cluster.
%     (d) §1.1 positioning paragraph added: explicit
%         distinction of the construction from operator-
%         theoretic programmes (Connes, Berry-Keating,
%         de Branges, Burnol), and classical-dictionary
%         translation of "local curvature" as the truncated
%         prime-sum side of the Weil explicit formula together
%         with the archimedean term, citing Edwards as the
%         standard source.
%     (e) Selberg cited at the formal asymptotic
%         Z(g^*) ~ D log kappa log log kappa in Open
%         Problem 2; removes orphan-bibitem status.
%   - No changes to mathematical content, proofs, or numerical
%     values. All 20 mathematical statements of v5/v6 retained.
%
% v6 (April 2026)
%   - Series template v2.0 compliance sweep:
%     (a) Notation and Conventions section added between
%         Abstract and Introduction (template invariant).
%     (b) PDF Unicode hygiene: \usepackage{cmap} added before
%         \usepackage[T1]{fontenc}; \emergencystretch=1em set.
%     (c) Source hygiene: internal project codename removed
%         from the GitHub bibliography entry.
%     (d) Document structure aligned with template v2.0:
%         Introduction has three subsections (Background and
%         Motivation, Main Results, Reader Contract).
%         Open Problems section added with stable labels.
%         Non-Circularity is an own section after Open Problems.
%         Acknowledgments and Call for Feedback as unnumbered
%         sections. \newpage before References.
%     (e) Bibliography reformatted to Concept-DOI style:
%         no version numbers in series bibitems; quotation-mark
%         titles for consistency across the series.
%     (f) Theorem environments numbered within section;
%         theorem titles use sentence-case "; proved" style.
%     (g) Footer carries motto "Pauca sed matura" exactly once.
%     (h) Change Log section removed from PDF body
%         (changelog kept in source comments only).
%   - No changes to mathematical content, proofs, or numerical
%     values. All 20 mathematical statements of v5 retained.
% ============================================================

\documentclass[11pt]{article}
\usepackage{cmap}
\usepackage[T1]{fontenc}
\usepackage[utf8]{inputenc}
\usepackage[a4paper,margin=1in]{geometry}
\usepackage{amsmath,amssymb,amsthm,mathtools}
\usepackage{xcolor}
\usepackage{hyperref}
\usepackage{enumitem}
\usepackage{setspace}
\usepackage{booktabs}
\usepackage{array}
\usepackage{graphicx}

\hypersetup{colorlinks=true,
  linkcolor=blue!50!black,
  urlcolor=blue!50!black,
  citecolor=blue!50!black}
\setlength{\parskip}{0.5em}
\setlength{\parindent}{0pt}
\onehalfspacing
\emergencystretch=1em

% --------------------------------------------------------
% Theorem environments (numbered within section)
% --------------------------------------------------------
\newtheorem{theorem}{Theorem}[section]
\newtheorem{proposition}[theorem]{Proposition}
\newtheorem{lemma}[theorem]{Lemma}
\newtheorem{corollary}[theorem]{Corollary}
\newtheorem{definition}[theorem]{Definition}
\newtheorem{assumption}[theorem]{Assumption}
\newtheorem{observation}[theorem]{Observation}
\newtheorem{conjecture}[theorem]{Conjecture}
\newtheorem*{remark}{Remark}
\newtheorem{openproblem}[theorem]{Open Problem}

% --------------------------------------------------------
% Series-wide macros
% --------------------------------------------------------
\newcommand{\Hlocal}{H_{\mathrm{local}}}
\newcommand{\Hren}{H_{\mathrm{local}}^{\mathrm{ren}}}
\newcommand{\Sad}{S_{\mathrm{ad}}}
\newcommand{\cren}{c_p^{\mathrm{ren}}}
\newcommand{\Zren}{Z_{\mathrm{ren}}}

% --------------------------------------------------------
% Title
% --------------------------------------------------------
\title{\vspace{-1.5em}\bfseries
From Local Curvature to the Weil Functional:\\[0.4em]
\large An Explicit Construction}
\author{Ulrich Tehrani}
\date{Research Note \textperiodcentered\ May 2026
      \textperiodcentered\ v10}

% ============================================================
\begin{document}
\maketitle
% ============================================================

\begin{abstract}
Building on the curvature decomposition established
in~\cite{Paper1}, we construct an explicit family of admissible
test functions $g^*_{\sigma,\varepsilon}\in\Sad$ for the Weil
explicit formula and identify the corresponding prime-side
contribution with the local curvature
$\Hlocal(\sigma,\kappa)$ studied in~\cite{Paper1}.
A renormalised family of prime weights
$\cren=f_p^{1/2}>0$ is introduced, and the associated diagonal
energy $D=\sum_p (\cren)^2\approx 9.471$ is shown to converge.
The renormalised local curvature is shown to be of lower order
than $(\log\kappa)^2$ and to exhibit a sawtooth structure governed,
on average, by the Mertens constant.
No hypothetical input is used: no Riemann Hypothesis, no
Hilbert--P\'olya postulate, no GUE conjecture.

\textbf{What is proved.}
The antisymmetrised Gaussian family $g^*_{\sigma,\varepsilon}$
belongs to the admissible class $\Sad$ used in
Lagarias~\cite{Lagarias} (Lemma~\ref{lem:admissible}).
The renormalised prime weight
$\cren := \bigl(V_p(\tfrac12)-\tfrac{4(\log p)^2}{p}\bigr)^{1/2}$
is strictly positive (Lemma~\ref{lem:fp-pos}), and the diagonal
energy series $D=\sum_p (\cren)^2$ converges absolutely
(Lemma~\ref{lem:D-conv}).

\textbf{What is conditional.}
The prime-side identity
\[
W(g^*_{\sigma,\varepsilon}*\check g^*_{\sigma,\varepsilon})
\;=\; Z(g^*_{\sigma,\varepsilon})\;-\;\Hlocal(\sigma,\kappa)
\;+\; R(\sigma,\varepsilon,\kappa)
\]
holds with $R(\sigma,\varepsilon,\kappa)\to 0$ as
$\varepsilon\to 0$ at fixed $(\sigma,\kappa)$
(Proposition~\ref{prop:prime-side}), conditional on the standard
remainder bookkeeping for the Weil explicit formula applied to
antisymmetric Schwartz test functions.
The remainder is not made uniform in $\kappa$.

\textbf{What is numerical.}
At reference parameters $(\varepsilon,N)=(0.05,100)$, the
diagonal energy evaluates to $D\approx 9.471$ on primes
below~$10^4$, with tail bound $<0.009$.
Direct grid evaluation of the renormalised quantities at
$\kappa\in\{10,20,29,50,100,200,500,1000\}$ shows that
$\Zren\ge\Hren$ holds at every tested point except $\kappa=20$,
where $\Zren-\Hren=-0.988$, with the ratio $\Zren/\Hren$ growing
from $1.5$ at $\kappa=29$ to over $10$ at $\kappa=1000$.
No claim is made about $\kappa$ outside the eight tested points.

\textbf{What remains open.}
Two questions separate the present construction from a full
prime-side derivation of Weil positivity for this family
(\S\,\ref{sec:open}).
Problem~\ref{op:spectral-ineq} (\emph{Spectral inequality})
asks whether $Z(g^*_{1/2,\varepsilon})\ge\Hlocal(\tfrac12,\kappa)$
holds unconditionally; an unconditional proof for this family
alone is necessary, but not sufficient, for the Riemann
Hypothesis in the absence of a density argument
in~$\Sad$.
Problem~\ref{op:offdiag} (\emph{Off-diagonal control}) asks for
an unconditional bound on the cross-terms
$\sum_\gamma\cos(\gamma\log(p/q))$ for $p\ne q$, which would
upgrade the formal asymptotic
$Z(g^*)\sim D\log\kappa\log\log\kappa$ to a proved statement.

\textbf{Structure.}
The formal core in \S\S\,\ref{sec:construct}--\ref{sec:scale} is
non-circular, with proved and conditional components explicitly
separated.
No use of the Riemann Hypothesis, the GUE conjecture, or the
Hilbert--P\'olya postulate is made anywhere in
\S\S\,\ref{sec:construct}--\ref{sec:scale}.
\end{abstract}

\vspace{2em}

% --------------------------------------------------------
% Notation and Conventions
% --------------------------------------------------------
\section*{Notation and Conventions}
\label{sec:notation}

\emph{Critical line.}
Throughout, the critical line means $\Re(s)=\tfrac12$.
The construction parameter $\sigma>0$ is evaluated at
$\sigma=\tfrac12$ in the renormalised analysis unless stated
otherwise.

\medskip
\emph{Global parameters.}
The reference values are $\kappa\ge 2$ (prime cutoff,
varied across the comparison grid),
$\varepsilon=0.05$ (Gaussian smoothing parameter), and
$N=100$ (number of Riemann ordinates $\gamma_k$ used in the
zero-sum evaluations of~\S\,\ref{sec:numerical}).
Values are fixed within each statement and explicitly varied
otherwise.

\medskip
\emph{Main objects.}
The local curvature, imported from~\cite{Paper1}, is
\[
\Hlocal(\sigma,\kappa) \;=\; \psi_1(\sigma)
\;+\; \sum_{p\le\kappa} V_p(\sigma),
\qquad
V_p(\sigma) \;=\; 4(\log p)^2\,
\frac{p^{-2\sigma}}{(1-p^{-2\sigma})^2}.
\]
The Gaussian bump is $\varphi_\varepsilon(x)=
(\varepsilon\sqrt{2\pi})^{-1} e^{-x^2/(2\varepsilon^2)}$.
The raw weight at level $\sigma$ is
$c_{p,\sigma} = V_p(\sigma)^{1/2}\cdot p^{1/4}/(\log p)^{1/2}$,
and the antisymmetrised admissible test function
$g^*_{\sigma,\varepsilon}\in\Sad$ is constructed
in~\S\,\ref{sec:construct}.
The Weil quadratic functional is denoted $W$, and
$Z(g)=\sum_\rho |\widehat g(\rho)|^2$ is the corresponding
zero-side sum over the non-trivial zeros of $\zeta$.
At $\sigma=\tfrac12$ the renormalised prime weight is
$\cren = f_p^{1/2}$, where
$f_p = V_p(\tfrac12)-4(\log p)^2/p$, and the diagonal
energy is $D=\sum_p (\cren)^2$.
The renormalised local curvature is
$\Hren(\tfrac12,\kappa)=\Hlocal(\tfrac12,\kappa)-2(\log\kappa)^2$.

\medskip
\emph{Not to be confused with.}
The diagonal energy $D\approx 9.471$ defined here, characteristic
of the renormalised prime-side construction of this paper, must
not be confused with the trace-formula diagonal $D_{\mathrm{SEL}}$
appearing in later work in the series, which is defined in a
different operator context and takes a different numerical value.
The renormalised weight $\cren=f_p^{1/2}$ used in the Weil
functional differs in normalisation from the unrenormalised
weight $c_{p,\sigma}$ used in the construction of
$g^*_{\sigma,\varepsilon}$.

\newpage

% --------------------------------------------------------
% §1 Introduction
% --------------------------------------------------------
\section{Introduction}
\label{sec:intro}

\subsection{Background and Motivation}

This paper is the second in a series developing an
operator-theoretic framework for studying the distribution of
the non-trivial zeros of the Riemann zeta function $\zeta(s)$.
The central organising principle of the series is that prime
numbers mediate coupling between zero ordinates, and that
explicit finite-cutoff constructions on the prime side carry
arithmetic information accessible by classical means.

\textbf{Paper~1}~\cite{Paper1} introduced the curvature field
\[
H_\xi(\sigma,t) \;:=\; \frac{d^2}{d\sigma^2}\,
\log|\xi(\sigma+it)|^2
\;=\; \Hlocal(\sigma,\kappa) \;+\; H_{\mathrm{dual}}(\sigma,t,\kappa),
\]
where $H_{\mathrm{dual}}=H_\xi-\Hlocal$ is the spectral
remainder encoding zero contributions, and where the truncated
local curvature
\[
\Hlocal(\sigma,\kappa) \;=\; \psi_1(\sigma)
\;+\; \sum_{p\le\kappa} V_p(\sigma)
\]
contains the archimedean term and the finite-place
contributions.
The prime-side pieces satisfy $V_p(\sigma)\ge 0$, and on the
critical line one has $\Hlocal(\tfrac12,\kappa)\asymp
(\log\kappa)^2$.

\textbf{The present paper} addresses a concrete part of the
central question raised in~\cite[\S\,7]{Paper1}: whether the
local positivity and critical divergence of $\Hlocal$ can be
organised into an explicit admissible test-function construction
whose prime-side term in the Weil functional matches the local
curvature.
We answer this in the affirmative on the prime side by
constructing an explicit family $g^*_{\sigma,\varepsilon}\in\Sad$
and identifying the corresponding prime-side contribution with
$\Hlocal(\sigma,\kappa)$ up to a smoothing remainder that
vanishes as $\varepsilon\to 0$.
The remaining open step --- the spectral (zero-side) inequality
on the critical line --- is then isolated as a single
RH-strength question in~\S\,\ref{sec:open}.

The construction $g^*_{\sigma,\varepsilon}\in\Sad$ developed here
should be distinguished from the established operator-theoretic
programmes addressing the Riemann Hypothesis.
It is not a Hilbert--P\'olya spectral realisation in the sense of
Connes~\cite{Connes} or Berry--Keating~\cite{BerryKeating}, nor a
de~Branges structure-function construction~\cite{deBranges}; it
is an explicit antisymmetrised Gaussian smearing of prime delta
masses with arithmetic weights, closest in spirit to Hilbert-space
reformulations of the explicit formula such as those of
Burnol~\cite{Burnol}.
The local curvature $\Hlocal(\sigma,\kappa)$ is, in the classical
dictionary of analytic number theory~\cite{Edwards}, the truncated
prime-sum side of the Weil explicit formula together with the
archimedean term $\psi_1(\sigma)$; the curvature label reflects the
second-derivative origin of $V_p(\sigma)$ in $\log|\xi|^2$
identified in~\cite{Paper1} and is not used in any further
geometric sense.
The terms ``prime--zero coupling'' and ``curvature field'' used
informally in the series are translations of the explicit-formula
duality and of $\partial_\sigma^2\log|\xi|^2$ respectively, not
new mathematical structures.

The construction is set within the Lagarias formulation of Weil
positivity~\cite{Lagarias}, which characterises the Riemann
Hypothesis as positivity of the Weil functional $W(f*\check f)$
on a suitable class of test functions, building on the Weil
explicit formula~\cite{Weil} and the Li positivity
criterion~\cite{Li}.
The renormalised diagonal scale $D=\sum_p(\cren)^2$ encountered
in~\S\,\ref{sec:renorm} is closely related, via the prime number
theorem and Mertens' classical estimates~\cite{Mertens}, to the
average behaviour of the local curvature; this connection is
made precise in~\S\,\ref{sec:scale}.

\subsection{Main Results}

The paper contains three main components.

The first is the explicit construction of the test family
$g^*_{\sigma,\varepsilon}$ and the proof that it lies in the
admissible class $\Sad$ used by Lagarias~\cite{Lagarias}.
The construction proceeds in two steps: a Gaussian smearing of
delta masses at $\log p$ with weights $c_{p,\sigma}$, followed
by antisymmetrisation enforcing the vanishing condition at the
origin (\S\,\ref{sec:construct}).
The admissibility lemma (Lemma~\ref{lem:admissible}) is a
direct verification.

The second is the prime-side identity
(Proposition~\ref{prop:prime-side}), established conditionally
on the standard remainder bookkeeping for the Weil explicit
formula applied to antisymmetric Schwartz test functions:
\[
W(g^*_{\sigma,\varepsilon}*\check g^*_{\sigma,\varepsilon})
\;=\; Z(g^*_{\sigma,\varepsilon})\;-\;\Hlocal(\sigma,\kappa)
\;+\; R(\sigma,\varepsilon,\kappa),
\qquad R\to 0\text{ as }\varepsilon\to 0
\text{ at fixed }(\sigma,\kappa).
\]
The remainder $R$ collects the smoothing-induced off-diagonal
cross-terms ($p\ne q$) together with the higher prime-power
contributions; it is not made uniform in $\kappa$ in the
present note.
The constant term in the explicit formula vanishes by the
antisymmetry of $g^*$.

The third is the renormalised analysis at the critical line
(\S\S\,\ref{sec:renorm}--\ref{sec:scale}).
At $\sigma=\tfrac12$ the auxiliary quantity
$f_p=V_p(\tfrac12)-4(\log p)^2/p$ is shown to satisfy the
closed form
\[
f_p \;=\; \frac{4(\log p)^2(2p-1)}{p(p-1)^2} \;>\; 0
\qquad (p\ge 2)
\]
(Lemma~\ref{lem:fp-pos}), so that $\cren=f_p^{1/2}$ is real
and strictly positive.
The diagonal energy $D=\sum_p(\cren)^2$ converges absolutely by
comparison with $\sum_n (\log n)^2/n^2$ (Lemma~\ref{lem:D-conv})
and evaluates to $D\approx 9.471$.
The renormalised local curvature
$\Hren(\tfrac12,\kappa)=\Hlocal(\tfrac12,\kappa)-2(\log\kappa)^2$
is of lower order than $(\log\kappa)^2$ and exhibits a sawtooth
structure: a positive jump $f_p$ at each new prime
$p\le\kappa$, decreasing continuously between primes.
The Mertens constant $\gamma_M$ governs the average drift of
this sawtooth via the classical estimate
$\sum_{p\le\kappa}(\log p)/p=\log\kappa+O(1)$.

The numerical comparison of~\S\,\ref{sec:numerical} establishes
that $\Zren\ge\Hren$ holds at every point of the eight-point
test grid except $\kappa=20$, with the ratio growing from $1.5$
at $\kappa=29$ to over $10$ at $\kappa=1000$.
Two open problems remain:
the unconditional spectral inequality
$Z(g^*_{1/2,\varepsilon})\ge\Hlocal(\tfrac12,\kappa)$
(Problem~\ref{op:spectral-ineq}) and
the off-diagonal control of the cross-terms in the zero sum
(Problem~\ref{op:offdiag}).

\subsection{Reader Contract}

This paper \emph{proves}, by direct construction and classical
identities, the admissibility of the test family
(Lemma~\ref{lem:admissible}), the strict positivity
of the renormalised weight $\cren=f_p^{1/2}>0$
(Lemma~\ref{lem:fp-pos}), and the absolute convergence of the
diagonal energy series $D=\sum_p(\cren)^2$
(Lemma~\ref{lem:D-conv}).
These are unconditional mathematical statements.

This paper \emph{establishes conditionally}, conditional on the
standard remainder bookkeeping for the Weil explicit formula
applied to antisymmetric Schwartz test functions, the prime-side
identity for the Weil functional at fixed $(\sigma,\kappa)$ as
$\varepsilon\to 0$ (Proposition~\ref{prop:prime-side}).
The remainder is not made uniform in $\kappa$ in the present
note; the off-diagonal control needed for a uniform statement is
the subject of Problem~\ref{op:offdiag}.

This paper \emph{establishes numerically}, at reference
parameters, the diagonal energy value $D\approx 9.471$
(Observation~\ref{obs:D-value}) and the finite-grid comparison
$\Zren\ge\Hren$ on the eight-point grid
$\kappa\in\{10,20,29,50,100,200,500,1000\}$ at every tested
point except $\kappa=20$ (Observation~\ref{obs:Z-vs-H}).
Both numerical statements are finite computations from the
stated parameters and are labelled [numerical] in their theorem
headings.

This paper \emph{leaves open} two well-posed questions named in
\S\,\ref{sec:open}:
Problem~\ref{op:spectral-ineq} (Spectral inequality) and
Problem~\ref{op:offdiag} (Off-diagonal control).
Both are formulated entirely within classical analytic number
theory.

% ============================================================
% §2 Construction
% ============================================================
\section{Construction of the admissible test family}
\label{sec:construct}

Recall from~\cite[Eq.~(3)]{Paper1} that
\[
V_p(\sigma)
\;=\; 4(\log p)^2\,\frac{p^{-2\sigma}}{(1-p^{-2\sigma})^2}.
\]
Define the raw weight
\[
c_{p,\sigma}
\;:=\; V_p(\sigma)^{1/2}
\cdot \frac{p^{1/4}}{(\log p)^{1/2}},
\]
and let
\[
\varphi_\varepsilon(x)
\;=\; \frac{1}{\varepsilon\sqrt{2\pi}}\,e^{-x^2/(2\varepsilon^2)}
\]
be the standard Gaussian bump.
For fixed $\kappa$ set
\[
\widehat f_{\sigma,\varepsilon}(x)
\;=\; \sum_{p\le\kappa} c_{p,\sigma}\,
       \varphi_\varepsilon(x-\log p).
\]
The admissible test function is obtained by antisymmetrisation:
\begin{equation}\label{eq:gstar}
\widehat g^*_{\sigma,\varepsilon}(x)
\;:=\; \widehat f_{\sigma,\varepsilon}(x)
       \;-\; \widehat f_{\sigma,\varepsilon}(-x).
\end{equation}

\begin{lemma}[Admissibility; proved]\label{lem:admissible}
For every $\sigma>0$ and $\varepsilon>0$, the function
$g^*_{\sigma,\varepsilon}$ belongs to the admissible class $\Sad$
used by Lagarias~\cite{Lagarias} in the formulation of Weil
positivity.
\end{lemma}

\begin{proof}
The Gaussian bumps $\varphi_\varepsilon$ are Schwartz, and a
finite linear combination of translated Schwartz functions is
again Schwartz; antisymmetrisation preserves the Schwartz class.
Hence $g^*_{\sigma,\varepsilon}$ and
$\widehat g^*_{\sigma,\varepsilon}$ are rapidly decreasing.
By antisymmetry, $\widehat g^*_{\sigma,\varepsilon}(0)=0$ and
$g^*_{\sigma,\varepsilon}(0)=0$, providing the vanishing
conditions at the origin required for admissibility.
The Gaussian decay ensures absolute convergence of the
prime-side and zero-side expressions appearing in the Weil
functional for this family.
\end{proof}

% ============================================================
% §3 Prime-side identity
% ============================================================
\section{The prime-side identity}
\label{sec:prime-side}

\begin{proposition}[Prime-side identity; conditional]\label{prop:prime-side}
For the family $g^*_{\sigma,\varepsilon}$ constructed
in~\S\,\ref{sec:construct},
\begin{equation}\label{eq:weil-decomp}
W(g^*_{\sigma,\varepsilon}*\check g^*_{\sigma,\varepsilon})
\;=\; Z(g^*_{\sigma,\varepsilon})\;-\;\Hlocal(\sigma,\kappa)
\;+\; R(\sigma,\varepsilon,\kappa),
\end{equation}
where $Z(g^*_{\sigma,\varepsilon})=\sum_\rho
|\widehat g^*_{\sigma,\varepsilon}(\rho)|^2$ runs over the
non-trivial zeros $\rho$ of $\zeta(s)$, and the remainder
$R(\sigma,\varepsilon,\kappa)$ tends to zero as $\varepsilon\to 0$
at fixed $(\sigma,\kappa)$.
\end{proposition}

\begin{proof}[Sketch]
The argument rests on the standard remainder bookkeeping for
the Weil explicit formula applied to antisymmetric Schwartz test
functions at fixed $(\sigma,\kappa)$ as $\varepsilon\to 0$, which
is the assumption under which the proposition is stated.

The Weil quadratic functional decomposes, in the Lagarias
normalisation~\cite{Lagarias}, into four pieces: a constant
term, an archimedean integral, a prime sum, and a zero sum.
The constant term, proportional to
$\widehat g^*_{\sigma,\varepsilon}(0)^2$, vanishes by
$\widehat g^*_{\sigma,\varepsilon}(0)=0$
(Lemma~\ref{lem:admissible}).
The archimedean integral, an integral against the
$\Gamma$-factor of $\xi$, contributes $\psi_1(\sigma)$ in the
limit $\varepsilon\to 0$ at fixed $\sigma$.
The zero-side sum is $Z(g^*_{\sigma,\varepsilon})$ by definition.
The prime-side quadratic form, expanded against the
double sum in the convolution
$g^*_{\sigma,\varepsilon}*\check g^*_{\sigma,\varepsilon}$,
contributes
\[
\sum_{p,q\le\kappa} c_{p,\sigma} c_{q,\sigma}
\bigl[\varphi_\varepsilon*\check\varphi_\varepsilon\bigr](\log p-\log q)
\cdot p^{-1/2} q^{-1/2}
\]
together with prime-power terms of higher arithmetic order.
By construction of the weights $c_{p,\sigma}$, the diagonal
$p=q$ contribution converges to
$\sum_{p\le\kappa}V_p(\sigma)$ as $\varepsilon\to 0$.
The limit of the prime-side together with the archimedean
integral therefore equals $\Hlocal(\sigma,\kappa)$.
The remainder $R(\sigma,\varepsilon,\kappa)$
collects the smoothing-induced off-diagonal cross-terms
($p\ne q$) together with the higher prime-power contributions;
both vanish as $\varepsilon\to 0$ at fixed $(\sigma,\kappa)$ by
Gaussian decay of $\varphi_\varepsilon*\check\varphi_\varepsilon$
at non-zero argument.
The remainder is not made uniform in $\kappa$ in the present
note; the off-diagonal control needed for a uniform statement
is the subject of Problem~\ref{op:offdiag} of~\S\,\ref{sec:open}.
\end{proof}

\begin{remark}
Equation~\eqref{eq:weil-decomp} is the bridge between the
curvature field of~\cite{Paper1} and the Weil-positivity
framework of~\cite{Lagarias}.
The vanishing of the constant term, by condition (S2) on the
admissible class, is the structural reason that the right-hand
side of~\eqref{eq:weil-decomp} carries no free additive term:
the archimedean integral does not vanish but contributes
$\psi_1(\sigma)$, which is then absorbed into
$\Hlocal(\sigma,\kappa)$.
\end{remark}

% ============================================================
% §4 Renormalised weights and convergence
% ============================================================
\section{Renormalised weights and convergence}
\label{sec:renorm}

At $\sigma=\tfrac12$ define
\[
f_p \;:=\; V_p(\tfrac12) \;-\; \frac{4(\log p)^2}{p}.
\]
A direct calculation gives
\begin{equation}\label{eq:fp}
f_p \;=\; \frac{4(\log p)^2(2p-1)}{p(p-1)^2} \;>\; 0
\qquad (p\ge 2).
\end{equation}

\begin{lemma}[Positivity of the renormalised weight; proved]\label{lem:fp-pos}
For every prime $p\ge 2$, the quantity $f_p$
in~\eqref{eq:fp} is strictly positive.
\end{lemma}

\begin{proof}
The numerator $2p-1>0$ and the denominator $p(p-1)^2>0$ for
every prime $p\ge 2$.
\end{proof}

The renormalised weight is then simply
\begin{equation}\label{eq:cren}
\cren \;:=\; f_p^{1/2} \;>\; 0,
\end{equation}
so that $(\cren)^2 = f_p$ exactly.
Numerical values for small primes:

\begin{center}
\begin{tabular}{cccc}
\toprule
$p$ & $V_p(\tfrac12)$ & $4(\log p)^2/p$ & $f_p = (\cren)^2$ \\
\midrule
2  & 3.844 & 0.961 & 2.883 \\
3  & 3.621 & 1.609 & 2.012 \\
5  & 3.238 & 2.072 & 1.166 \\
7  & 2.945 & 2.164 & 0.781 \\
11 & 2.530 & 2.091 & 0.439 \\
\bottomrule
\end{tabular}
\end{center}

\begin{lemma}[Convergence of the diagonal energy; proved]\label{lem:D-conv}
The series
\begin{equation}\label{eq:D}
D \;:=\; \sum_p (\cren)^2
\;=\; \sum_p \frac{4(\log p)^2(2p-1)}{p(p-1)^2}
\end{equation}
converges absolutely.
\end{lemma}

\begin{proof}
For large $p$, $(\cren)^2 \sim 8(\log p)^2/p^2$.
Since $\sum_n (\log n)^2/n^2 < \infty$, absolute convergence
follows by comparison with this dominating series.
\end{proof}

\begin{observation}[Diagonal energy value; numerical]\label{obs:D-value}
Numerically, $D\approx 9.471$ when summed over primes below
$10^4$, with tail bound below $0.009$.
\end{observation}

% ============================================================
% §5 Renormalised scale and Mertens constant
% ============================================================
\section{The renormalised scale and the Mertens constant}
\label{sec:scale}

By the prime number theorem,
\[
\Hlocal(\tfrac12,\kappa) \;\sim\; 2(\log\kappa)^2
\qquad (\kappa\to\infty),
\]
so the renormalised remainder
\[
\Hren(\tfrac12,\kappa) \;:=\;
\Hlocal(\tfrac12,\kappa) \;-\; 2(\log\kappa)^2
\]
is of lower order than $(\log\kappa)^2$.
Numerically, $\Hren$ exhibits a sawtooth pattern: at each new
prime $p\le\kappa$ it jumps upward by $f_p>0$, and decreases
continuously between primes as $2(\log\kappa)^2$ grows.
This is the direct analogue of the prime-counting function
$\pi(x)$, which increases by~$1$ at each prime and is constant
between consecutive primes.
The Mertens constant $\gamma_M$ governs the average drift of
this sawtooth via the classical estimate
\[
\sum_{p\le\kappa}\frac{\log p}{p} \;=\; \log\kappa + O(1)
\qquad (\kappa\to\infty),
\]
connecting the renormalised scale to classical prime-sum
estimates~\cite{Mertens}.

\begin{remark}
The Mertens constant
\[
\gamma_M \;=\; \gamma + \sum_p\Bigl[\log\bigl(1-\tfrac1p\bigr)+\tfrac1p\Bigr]
\;\approx\; 0.2615,
\]
where $\gamma\approx 0.5772$ is the Euler--Mascheroni constant,
appears in prime-sum estimates that control the average
behaviour of $\Hren$.
The same constant $\gamma$ enters the first Li
coefficient~\cite{Li}
\[
\lambda_1 \;=\; 1 + \tfrac{\gamma}{2} - \log 2
- \tfrac{\log\pi}{2} \;\approx\; 0.0231,
\]
reflecting the universal role of $\gamma$ in the transition
from discrete prime sums to continuous logarithmic integrals.
\end{remark}

\begin{remark}
The renormalised construction clarifies the scale of the
prime-side contribution but does not by itself resolve the
spectral inequality of Problem~\ref{op:spectral-ineq}.
\end{remark}

% ============================================================
% §6 Numerical comparison
% ============================================================
\section{Numerical comparison of zero and prime sides}
\label{sec:numerical}

We compare the renormalised zero sum $\Zren$ with the
renormalised local curvature $\Hren$ on a finite grid in
$\kappa$.
The smoothing parameter $\varepsilon=0.05$ is required for
meaningful numerical experiments with the zero sum: for
$\varepsilon=0.5$ one has $e^{-\varepsilon^2\gamma_1^2}\approx
10^{-22}$, effectively suppressing all zero contributions, so
that no observable signal remains.
With $\cren=f_p^{1/2}$, $\varepsilon=0.05$, and $N=100$ zeros:

\begin{center}
\begin{tabular}{ccccc}
\toprule
$\kappa$ & $\Zren$ & $\Hren$
         & $\Zren-\Hren$ & $\Zren\ge\Hren$ \\
\midrule
10   & 4.888 & 3.044 & $+1.844$ & $\checkmark$ \\
20   & 3.782 & 4.770 & $-0.988$ & $\times$ \\
29   & 5.520 & 3.588 & $+1.931$ & $\checkmark$ \\
50   & 5.693 & 2.881 & $+2.812$ & $\checkmark$ \\
100  & 5.311 & 1.562 & $+3.748$ & $\checkmark$ \\
200  & 5.102 & 2.355 & $+2.747$ & $\checkmark$ \\
500  & 5.156 & 1.101 & $+4.056$ & $\checkmark$ \\
1000 & 4.289 & 0.426 & $+3.863$ & $\checkmark$ \\
\bottomrule
\end{tabular}
\end{center}

\begin{observation}[Finite-grid stability of the comparison; numerical]\label{obs:Z-vs-H}
On the eight-point grid
$\kappa\in\{10,20,29,50,100,200,500,1000\}$, the inequality
$\Zren\ge\Hren$ holds at every tested point except $\kappa=20$,
where $\Zren-\Hren=-0.988$.
The ratio $\Zren/\Hren$ grows from $1.5$ at $\kappa=29$ to
over $10$ at $\kappa=1000$.
No claim is made about $\kappa$ outside the eight tested points.
\end{observation}

The single grid violation $\Zren<\Hren$ at $\kappa=20$ illustrates
the sawtooth behaviour of $\Hren$: directly after a prime jump
$\Hren$ rises by $f_p$, then decreases continuously until the
next prime, exactly as $\pi(x)-\mathrm{Li}(x)$ fluctuates near
small primes before stabilising asymptotically.
The pattern observed on the grid for $\kappa\ge 29$ is consistent
with the spectral inequality of Problem~\ref{op:spectral-ineq},
but does not prove it: the underlying object is the unconditional
positivity of the zero-side quadratic form for an admissible
test function in $\Sad$, which is not implied by finite-grid
evaluations.

\begin{remark}[The constant $D/2\pi$]
The formal asymptotic
$Z(g^*)\sim D\log\kappa\log\log\kappa$, derivable under
unconditional control of the off-diagonal cross-terms
(Problem~\ref{op:offdiag}), yields a natural bridge constant
\[
\frac{D}{2\pi} \;\approx\;
\frac{9.471}{6.283} \;\approx\; 1.507,
\]
measuring the ratio of the prime-side diagonal energy $D$ to
the archimedean normalisation $2\pi$ of the zero-counting
density $N(T)\sim T\log T/(2\pi)$.
This constant is of a different nature from the pair-correlation
statistics studied in the GUE conjecture~\cite{Montgomery}: it
depends only on the total density of zeros via the classical
formula for $N(T)$, which is unconditional and requires neither
GUE nor RH as input.
Whether $D/2\pi$ admits a representation in terms of classical
constants is left open.
\end{remark}

All numerical claims of this section are reproducible from the
verification scripts in the GitHub repository~\cite{Code}.

% ============================================================
% §7 Open Problems
% ============================================================
\section{Open Problems}
\label{sec:open}

The present work reduces the question of prime-side derivation
of Weil positivity for the family $g^*_{\sigma,\varepsilon}$ to
two well-posed subproblems.
We name each and indicate the tools we expect to be relevant.

\begin{openproblem}[Spectral inequality on the critical line]\label{op:spectral-ineq}
Prove that for every cutoff $\kappa$ and every sufficiently
small $\varepsilon>0$,
\begin{equation}\label{eq:openineq}
Z(g^*_{1/2,\varepsilon}) \;\ge\; \Hlocal(\tfrac12,\kappa).
\end{equation}
This inequality is of RH strength.
By Lagarias~\cite{Lagarias}, the equivalence
$\mathrm{RH}\iff W(f*\check f)\ge 0$ holds for all
$f\in\Sad$.
A positive answer to~\eqref{eq:openineq} establishes
$W(g^*_{1/2,\varepsilon}*\check g^*_{1/2,\varepsilon})\ge 0$ for
the specific family constructed in~\S\,\ref{sec:construct},
which is necessary, but not sufficient, for RH in the absence
of a density argument for this family in $\Sad$.
The expected tools are: explicit lower bounds on the zero-side
quadratic form by the unconditional zero-counting density;
control of the cross-terms in the zero sum
(Problem~\ref{op:offdiag}).
On the tested grid the inequality $\Zren\ge\Hren$ holds at every
tested point except $\kappa=20$ (Observation~\ref{obs:Z-vs-H}),
consistent with~\eqref{eq:openineq} but not establishing it
analytically.
\end{openproblem}

\begin{openproblem}[Off-diagonal control of the zero sum]\label{op:offdiag}
Prove unconditionally that the off-diagonal contribution
\[
\sum_{p\ne q}^{p,q\le\kappa} \cren \, c_q^{\mathrm{ren}}
\sum_{\rho} \widehat g^*_{1/2,\varepsilon}(\rho)\,
            \overline{\widehat g^*_{1/2,\varepsilon}(\rho)}
\,\Big|_{p\ne q\text{ part}}
\]
is asymptotically negligible compared with the diagonal
contribution
\[
\sum_p (\cren)^2 \cdot I(\varepsilon)
\;\sim\; D\,\log\kappa,
\]
so that, combined with the unconditional density
$N(T)\sim T\log T/(2\pi)$, the formal asymptotic
$Z(g^*)\sim D\log\kappa\log\log\kappa$, expected on the
Selberg-orthogonality heuristic for independent zero
contributions~\cite{Selberg}, becomes a proved
statement.
The diagonal part is unconditional: it follows from the prime
number theorem and the absolute convergence
of~\eqref{eq:D}.
The expected tools are: analytic continuation of
$\sum_\gamma\cos(\gamma\log(p/q))$ for $p\ne q$ via classical
explicit-formula techniques; cancellation by oscillation, with
the $(\log\kappa)^{-1}$-decay of the diagonal entering as a
size constraint.
This problem is formulated entirely within classical analytic
number theory.
\end{openproblem}

% ============================================================
% §8 Non-Circularity
% ============================================================
\section{Non-Circularity}
\label{sec:nocirc}

We document explicitly that the results of this paper do not
rely on the objects they are directed toward.

\textbf{No RH as input.}
At no point is the Riemann Hypothesis assumed.
The non-trivial zeros $\rho$ enter the construction only through
the classical zero-counting density $N(T)\sim T\log T/(2\pi)$,
which is unconditional, and through finite Gaussian-smoothed
sums for which the symmetric pairing $\rho\leftrightarrow\bar\rho$
is used without further hypothesis.
The reference point $\sigma=\tfrac12$ is chosen as the object
of study, not assumed to be the unique zero location.

\textbf{No GUE, Montgomery, or Hilbert--P\'olya hypothesis.}
The construction of $g^*_{\sigma,\varepsilon}$ uses only
Schwartz-class smoothing and antisymmetrisation; it makes no
reference to random-matrix statistics, pair-correlation
conjectures, or any postulated spectral realisation of the
zeros.
The bridge constant $D/2\pi$ identified in~\S\,\ref{sec:numerical}
depends only on the total density of zeros via the classical
formula for $N(T)$ and is structurally independent of GUE.

\textbf{The Weil explicit formula is classical.}
Proposition~\ref{prop:prime-side} applies the Weil explicit
formula~\cite{Weil} in the Lagarias admissible
formulation~\cite{Lagarias}; both are valid for every
sufficiently regular test function and require no positivity
hypothesis beyond the classical ones used in the derivation of
the explicit formula itself.

\textbf{Numerical results are clearly marked.}
The diagonal energy value $D\approx 9.471$ and the
comparison-grid statements of~\S\,\ref{sec:numerical} rely on
direct computation; each statement is labelled [numerical] in
its heading.

\textbf{Open problems are named, not used as hypotheses.}
Problems~\ref{op:spectral-ineq} and~\ref{op:offdiag} are stated
openly in~\S\,\ref{sec:open}.
No claim of this paper is conditional on a positive answer to
either problem.

% ============================================================
% Acknowledgments
% ============================================================
\section*{Acknowledgments}

The author thanks Jean-Fran\c{c}ois Burnol for confirming that
the curvature decomposition of~\cite{Paper1} does not appear in
the literature known to him, and the readers who provided
feedback on earlier versions of this paper.

% ============================================================
% Call for Feedback
% ============================================================
\section*{Call for Feedback}

This paper is part of an ongoing open research initiative
toward the Riemann Hypothesis.
All papers in the series are freely available on Zenodo under
CC~BY~4.0, and the accompanying verification code is hosted on
GitHub.

The author welcomes critical feedback, corrections, and
collaboration:
{\sloppy
Zenodo (\url{https://zenodo.org/search?q=Ulrich+Tehrani}) and
GitHub (\url{https://github.com/utehrani/analysislab-nt}).
\par}

Topics of particular interest include: analytical approaches to
the spectral inequality of Problem~\ref{op:spectral-ineq};
unconditional control of the off-diagonal cross-terms of
Problem~\ref{op:offdiag}; and the relation between the bridge
constant $D/2\pi$ identified in~\S\,\ref{sec:numerical} and
classical analytic constants attached to the zero-counting
density.

We particularly welcome expert comment on the following
question: for the Gaussian-smoothed admissible family
$g^*_{\sigma,\varepsilon}$ constructed in~\S\,\ref{sec:construct},
is the comparison $Z(g^*)\sim D\log\kappa\log\log\kappa$, with
$D$ as in~\eqref{eq:D}, accessible by classical sieve or
explicit-formula techniques without recourse to random-matrix
statistics?

% ============================================================
% References
% ============================================================
\newpage
\begin{thebibliography}{99}

\bibitem{Paper1}
U.~Tehrani,
``A Curvature Decomposition of the Explicit Formula'',
Zenodo, \url{https://doi.org/10.5281/zenodo.19025598}, 2026.

\bibitem{Code}
U.~Tehrani,
\emph{Verification scripts for the prime--zero coupling series},
GitHub, \url{https://github.com/utehrani/analysislab-nt}.

\bibitem{Lagarias}
J.C.~Lagarias,
``On a positivity property of the Riemann $\xi$-function'',
\emph{Acta Arithmetica} \textbf{89} (1999), 217--234.

\bibitem{Weil}
A.~Weil,
``Sur les `formules explicites' de la th\'eorie des nombres
premiers'',
\emph{Comm.\ S\'em.\ Math.\ Univ.\ Lund, Suppl.}\ (1952),
252--265.

\bibitem{Li}
X.-J.~Li,
``The positivity of a sequence of numbers and the Riemann
hypothesis'',
\emph{J.\ Number Theory} \textbf{65} (1997), 325--333.

\bibitem{Mertens}
F.~Mertens,
``Ein Beitrag zur analytischen Zahlentheorie'',
\emph{J.\ Reine Angew.\ Math.} \textbf{78} (1874), 46--62.

\bibitem{Montgomery}
H.L.~Montgomery,
``The pair correlation of zeros of the zeta function'',
in: \emph{Analytic Number Theory},
Proc.\ Sympos.\ Pure Math.\ \textbf{24},
Amer.\ Math.\ Soc., Providence, 1973, pp.~181--193.

\bibitem{Selberg}
A.~Selberg,
``Contributions to the theory of the Riemann zeta-function'',
\emph{Arch.\ Math.\ Naturvid.} \textbf{48}.5 (1946), 89--155.

\bibitem{Connes}
A.~Connes,
``Trace Formula in Noncommutative Geometry and the Zeros of the
Riemann Zeta Function'',
\emph{Selecta Math.\ (N.S.)} \textbf{5}.1 (1999), 29--106.

\bibitem{BerryKeating}
M.V.~Berry and J.P.~Keating,
``The Riemann Zeros and Eigenvalue Asymptotics'',
\emph{SIAM Review} \textbf{41}.2 (1999), 236--266.

\bibitem{deBranges}
L.~de~Branges,
``The Riemann Hypothesis for Hilbert Spaces of Entire Functions'',
\emph{Bull.\ Amer.\ Math.\ Soc.\ (N.S.)} \textbf{15}.1 (1986),
1--17.

\bibitem{Burnol}
J.-F.~Burnol,
``On Fourier and Zeta(s)'',
\emph{Forum Math.} \textbf{16}.6 (2004), 789--840.

\bibitem{Edwards}
H.M.~Edwards,
\emph{Riemann's Zeta Function},
Academic Press, New York, 1974.

\end{thebibliography}

% ============================================================
% Footer — normative format
% ============================================================
\enlargethispage{3\baselineskip}
\vspace*{\fill}
\begin{minipage}{\linewidth}
\rule{\linewidth}{0.4pt}\smallskip\\
\small\emph{Paper~2 v10 \textperiodcentered\ May~2026
\textperiodcentered\ Pauca sed matura}
\end{minipage}

\end{document}
